% reality/measurement_and_imperfect_data.tex

\chapter{Measurement and Imperfect Data}

\label{chap:measurement-and-imperfect-data}

\index{Observation!measurement data}
\index{Data quality!imperfect data}

Measurement provides observation evidence about realized infrastructure.
It does not replace constructive identity or realized-state modelling.

% ============================================================
\section{Observation Context}
% ============================================================

Measurement data in railway projects is heterogeneous in resolution,
accuracy, provenance, and coordinate context.

\begin{definition}[Observed data]
\index{Observed data}
Observed data denotes measured quantities acquired from sensors, surveys,
or legacy datasets in a specified observation context.
\end{definition}

\begin{principle}[Observation context principle]
\index{Observation context principle}
Observed data must be interpreted together with sensor, stationing,
coordinate, and topology context.
\end{principle}

% ============================================================
\section{Imperfect Data Characteristics}
% ============================================================

Observed data typically includes random noise, systematic bias,
incompleteness, and cross-source inconsistency.

\begin{definition}[Systematic error]
\index{Error!systematic}
A systematic error is a persistent observation deviation induced by
sensor, calibration, or transformation effects.
\end{definition}

\begin{definition}[Measurement uncertainty]
\index{Uncertainty!measurement}
Measurement uncertainty quantifies the interval or distribution of
plausible physical values compatible with an observation.
\end{definition}

% ============================================================
\section{Inference from Observation}
% ============================================================

Geometry and state are inferred from observation through reconstruction
operators and interpretation rules.

\begin{definition}[Discrete curvature estimate]
\index{Curvature!discrete estimate}
Given local directional samples, a discrete curvature estimate may be
written as
\[
\kappa_i \approx \frac{\theta_{i+1}-\theta_i}{\Delta s_i}.
\]
\end{definition}

Such estimates are sampling-dependent and uncertainty-sensitive;
filtering and regularization are therefore required in practice.

% ============================================================
\section{Observation and Runtime Reconstruction}
% ============================================================

Observation enters AIM through runtime reconstruction pipelines that
combine data with metric realization, stationing context, and topology
context.

Observed values are interpreted as evidence for realized state, not as
canonical object identity.

% ============================================================
\section{Engineering Use}
% ============================================================

In engineering workflows, imperfect data affects validation, maintenance
planning, and target--realized comparison. Robust workflows therefore
require provenance tracking, uncertainty-aware evaluation, and explicit
assumption management.
\index{Provenance!tracking}

% ============================================================
\section{Connection to AIM}
% ============================================================

\begin{summary}
In AIM, measurement and imperfect data form the observation layer of the
Reality part.

This layer supports realization-aware reconstruction and interpretation
while preserving the distinction between observed evidence and
engineering identity.
\end{summary}
